\Needspace{8\baselineskip}
\Article{212}
\textsc{Aufgabe.} \textit{Alle Darstellungen einer gegebenen Zahl $M$ durch eine gegebene Form $axx+2bxy+cyy$ mit der Determinante $hh$ finden.}

Die Lösung dieser Aufgabe lässt sich aus den Grundsätzen von Artikel~168 genauso ableiten, wie wir es oben (Artikel~180, 181, 205) für Formen mit negativer und mit positiver nichtquadratischer Determinante gezeigt haben. Da dies keine Schwierigkeit macht, wäre es überflüssig, es hier zu wiederholen. Es ist aber sinnvoll, die Lösung aus einem anderen Grundsatz abzuleiten, der dem vorliegenden Fall eigen ist.

Setzen wir wie in den Artikeln~206, 208
\[
\begin{gathered}
h-b:a=c:-(h+b)=\beta:\delta\\
\frac{h-b}{\beta}=\frac a\delta=f;\quad
\frac c\beta=\frac{-h-b}{\delta}=g
\end{gathered}
\]
so lässt sich leicht beweisen, dass die gegebene Form das Produkt der Faktoren $\delta x-\beta y$ und $fx-gy$ ist. Daher liefert jede Darstellung von $M$ durch die gegebene Form eine Zerlegung von $M$ in zwei Faktoren. Sind also alle Teiler von $M$ gleich $d,d',d''$ usw. (einschließlich $1$ und $M$, wobei jeder \textit{zweimal}, sowohl positiv als auch negativ, zu nehmen ist), so erhalten wir offenbar alle Darstellungen von $M$, indem wir nacheinander
\[
\begin{aligned}
\delta x-\beta y&=d,&\quad fx-gy&=\frac M d\\
\delta x-\beta y&=d',& fx-gy&=\frac M{d'}\quad\text{usw.}
\end{aligned}
\]
setzen, daraus die Werte von $x$, $y$ ermitteln und diejenigen Darstellungen verwerfen, bei denen $x$
\SourceLeaf{211}
oder $y$ gebrochene Werte annehmen. Aus den beiden ersten Gleichungen folgt offenbar
\[
x=\tfrac{\beta M-gdd}{(\beta f-\delta g)d}\quad\text{und}\quad
y=\tfrac{\delta M-fdd}{(\beta f-\delta g)d}
\]
Diese Werte sind immer \textit{bestimmt}, denn $\beta f-\delta g=2h$, und daher ist der Zähler sicher nicht $=0$. --- Auch die übrigen Aufgaben hätten sich aus demselben Grundsatz lösen lassen, nämlich dass jede Form mit quadratischer Determinante in zwei Faktoren zerlegt werden kann. Wir haben es aber vorgezogen, auch hier eine Methode zu verwenden, die der oben für Formen mit nichtquadratischer Determinante angegebenen entspricht.

\textit{Beispiel.} Gesucht sind alle Darstellungen von $12$ durch $3xx+4xy-7yy$. Diese Form zerfällt in die Faktoren $x-y$ und $3x+7y$. Alle Teiler von $12$ sind $\pm1,2,3,4,6,12$. Setzen wir $x-y=1$, $3x+7y=12$, so ergibt sich $x=\frac{19}{10}$, $y=\frac9{10}$; diese Werte müssen wir als gebrochene verwerfen. Ebenso ergeben die Teiler $-1,\pm3,\pm4,\pm6,\pm12$ unbrauchbare Werte. Der Teiler $+2$ liefert dagegen $x=2,y=0$, und $-2$ liefert $x=-2$, $y=0$. Außer diesen beiden gibt es also keine weiteren Darstellungen.

Diese Methode lässt sich nicht anwenden, wenn $M=0$ ist. In diesem Fall müssen offenbar alle Werte von $x,y$ entweder die Gleichung $\delta x-\beta y=0$ oder $fx-gy=0$ erfüllen. Alle Lösungen der ersten Gleichung sind in der Formel $x=\beta z$, $y=\delta z$ enthalten, wobei $z$ alle ganzen Zahlen durchläuft (denn $\beta$, $\delta$ sind wie vorausgesetzt teilerfremd). Setzen wir ebenso den größten gemeinsamen Teiler von $f,g$ $=m$, so werden alle Lösungen der zweiten Gleichung durch $x=\frac{gz}{m},y=\frac{hz}{m}$ dargestellt. Diese beiden allgemeinen Formeln umfassen daher alle Darstellungen von $M$ in diesem Fall.

\begin{center}
\(\begin{matrix} *&\qquad&*\\[-0.8ex]&*&\end{matrix}\)
\end{center}

Im Vorhergehenden haben wir alles behandelt, was zur Erkennung der Äquivalenz, zur Bestimmung aller Transformationen von Formen und zur Ermittlung aller Darstellungen gegebener Zahlen durch gegebene Formen gehört, sodass nichts mehr zu wünschen übrig zu bleiben scheint. Es bleibt daher nur noch zu zeigen, wie wir bei zwei gegebenen Formen, die wegen der \textit{Ungleichheit ihrer Determinanten} nicht äquivalent sein können, entscheiden, ob die eine unter der anderen enthalten ist, und in diesem Fall alle Transformationen der ersten in die zweite finden.

